% General Conclusion

\chapter*{General Conclusion}
\addcontentsline{toc}{chapter}{General Conclusion}

\section*{Summary of Achievements}
This project successfully designed and deployed a secure, privacy-respecting, \textbf{open-source alternative} to proprietary centralized social platforms. By utilizing a \textbf{Zero-Knowledge Architecture}, the platform addresses the surveillance risks associated with centralized data harvesting. 

Through E2EE text messaging and peer-to-peer (P2P) WebRTC voice calls, our implementation ensures that the signaling server acts strictly as a neutral relay, storing \textbf{no plaintext data}. While the platform currently utilizes server-managed keys for encryption to facilitate accessibility, the architecture is designed to be easily migrated to full client-side key custody. The development timeline was successfully executed across five distinct sprints under the Agile Scrum framework (facilitated by a dedicated Security Champion), providing a fully functional platform that preserves user digital sovereignty by default.

\section*{Cryptographic vs. UX Reflection}
A key technical challenge was balancing strong cryptographic constraints with a fluid, modern user experience. Secure architectures introduce performance challenges that required optimization:
\begin{itemize}
    \item \textbf{Symmetric Encryption Processing}: Performing client-side AES-256-GCM encryption and decryption on the browser thread requires careful optimization to prevent UI freezing. We addressed this by ensuring cryptographic operations do not block the main render loop.
    \item \textbf{Secure Search Operations}: Because the backend server does not store plaintext messages, global chat searches had to be re-engineered. This required implementing client-side indexing and searching over decrypted logs cached in local IndexedDB containers, ensuring search functionality without compromising data privacy on the server.
    \item \textbf{NAT Traversal Complexity}: Bypassing centralized media servers for P2P voice calling required configuring robust STUN/TURN connection fallback pipelines. This ensures connection reliability under restrictive corporate firewalls without compromising media confidentiality via DTLS/SRTP.
\end{itemize}
We resolved these bottlenecks by optimizing browser rendering cycles and indexing queries, demonstrating that robust security and high client responsiveness are not mutually exclusive.

\section*{Future Perspectives}
While the platform is fully completed, we have established a clear technology roadmap for future open-source developments to address current limitations and expand capabilities:
\begin{enumerate}
    \item \textbf{Independent Security Audit}: Preparing the codebase for an independent, third-party security audit of our SCRAM-SHA-256 authentication flows and AES-256-GCM encryption pipelines to ensure cryptographic integrity against advanced adversarial models.
    
    \item \textbf{Optimized P2P Mesh Video}: Expanding Sprint 4 audio pipelines to support multi-party P2P mesh video calling, utilizing SFU (Selective Forwarding Unit) architectures designed to protect peer IP privacy.
    
    \item \textbf{Contributor Documentation}: Publishing onboarding guides, security contribution protocols, and self-hosting documentation for the global open-source community.

    \item \textbf{Advanced Client-Side Search Indexing}: Although the architecture prevents server-side indexing of message content, the current client-side implementation using IndexedDB may face latency issues with large conversation histories. Future development will focus on integrating lightweight, browser-based Full-Text Search (FTS) engines to enable near-instant retrieval across extensive logs.

    \item \textbf{Migration to True Client-Side Key Management}: Currently, the platform utilizes a server-managed secret key for encryption. A critical future upgrade involves migrating to a decentralized key generation model where keys are generated and stored exclusively on the client device, achieving the strict Zero-Knowledge architecture where the server holds no decryption capability.

    \item \textbf{Cross-Device Synchronization}: Currently, the reliance on local IndexedDB containers binds data to a single device. A critical future upgrade involves implementing a secure multi-device synchronization protocol, allowing users to securely port their encrypted keys and message history between desktop and mobile clients without exposing sensitive data to the relay server.
\end{enumerate}
This senior thesis project at ITEAM successfully demonstrates that secure, decentralized communication is a viable, high-performance option for modern online communities.